% Chapter 3: Methodology (the paper plan's "Data and Labeling" and "Method").
% Project rules that hold throughout: time-based splits only, volatility models
% fitted on the training block only, fixed seeds (config.yaml, seed 42).
\label{ch:method}

\section{Introduction}
This chapter describes the data, the volatility models, the anomaly labels, the detectors
and the evaluation protocol. The study is quantitative and experimental. Two comparisons
drive it: a detector trained on raw returns against the same detector trained on
GARCH-standardized returns, and a detector trained on global market indices and applied to
the Dhaka Stock Exchange (DSE) against one trained on DSE data. Sections marked as planned
describe work that has not yet been run.

\section{Research Design}
\todo{Map each research question in Chapter 1 to the experiment that answers it.}

\section{Data Collection}

\subsection{Global Market Data}
Global index data come from Yahoo Finance through the \texttt{yfinance} Python package. The
first index used is the S\&P~500 (\texttt{\^{}GSPC}), daily, from January 2010 to December
2025. \todo{List the other global indices once chosen.}

The preliminary results in Chapter~4 were produced before Yahoo Finance could be reached
from the development environment. They use the copy of S\&P~500 daily prices that ships with
the Python \texttt{arch} package, which covers January 1999 to December 2018 and was
originally sourced from Yahoo Finance.

\subsection{Dhaka Stock Exchange Data}
The DSE publishes daily data for listed stocks in its day-end archive, with open, high,
low and close prices, the previous close (\texttt{ycp}), the number of trades, traded value
and volume. For long-listed stocks the archive goes back to about 1999. Index levels come
from the DSE market summary pages. Both sources are accessed through the open-source
\texttt{bdshare} package. A public dataset on Mendeley (DSE historical data, 1999 to April
2025, over 700 tickers) is used to cross-check the pull. It has no previous-close field, so
it cannot be used to detect corporate-action adjustments.

The market index is the DSE Broad Index (DSEX), which starts on 28 January 2013. Before
2013 the exchange published the DSE General Index (DGEN), which used a different method. The
two series are not spliced, so the DSE index sample begins in 2013.

Intraday DSE data are not freely available; the only historical source found is a paid
commercial feed. This thesis therefore uses daily data only.

\todo{Report the result of the DSE pull: number of DSEX rows, stocks retained and their
spans. The plan proceeds with daily data if DSEX has at least 3{,}000 rows from 2013 and at
least ten stocks have ten or more years with under 2\% missing days outside floor-price
periods. If the pull fails, the fallback market is India (NIFTY~50 and large-cap NSE
stocks through \texttt{yfinance}).}

\section{Data Cleaning and Preprocessing}
Returns are daily log returns in percent, $r_t = 100\ln(P_t/P_{t-1})$. Scaling by 100 keeps
the GARCH optimiser well conditioned.

DSE prices need four corrections before returns are computed.

\textbf{Floor prices.} The Bangladesh Securities and Exchange Commission (BSEC) imposed
minimum prices from about March 2020 to June 2021 and again from July 2022 to January 2024,
with partial lifting in December 2022 and reimposition in March 2023. During these periods
many stocks printed the same close with near-zero volume for months. Such flat runs are not
calm markets and must not be treated as normal behaviour. They are flagged as a separate
regime, and results are reported with and without them. The dates above come from press
reports at month level and will be checked against BSEC orders before the label rules are
fixed.

\textbf{Bonus shares and other corporate actions.} DSE prices are not adjusted, and
Bangladeshi firms issue bonus shares often, so raw closes fall sharply on record dates. The
DSE resets the previous close after a corporate action. A day on which \texttt{ycp} differs
from the prior day's close is treated as an adjustment day, and the earlier part of the
series is scaled by the ratio of the two. This rule is an inference about how the DSE sets
\texttt{ycp} and will be spot-checked against known bonus dates.

\textbf{Price limits.} Daily moves on the DSE are capped, at about $\pm 10\%$ for most of
the sample, so returns are truncated. Days that close at a limit are flagged rather than
treated as ordinary extreme returns.

\textbf{Calendar and liquidity.} DSEX trading dates define the calendar. Missing days and
zero-volume days are counted for each stock, and the final 10 to 20 stocks are chosen by
median daily turnover.

\section{Train, Validation and Test Splits}
All splits are by date, with no shuffling, so that no information from a later period is
used to fit or tune anything evaluated on an earlier one. For the S\&P~500 from
\texttt{yfinance}, the training block ends on 31 December 2018, validation covers 2019 to
2021, and the test block is everything from 2022 on. For the preliminary run on the bundled
1999 to 2018 series, training covers 1999 to 2012, validation 2013 to 2015 and test 2016 to
2018. \todo{DSE split dates.}

\section{Volatility Modelling and Standardization}
Each return is modelled as a constant mean plus an innovation, $r_t = \mu + \varepsilon_t$,
with $\varepsilon_t = \sigma_t e_t$ and $e_t$ drawn from a Student-$t$ distribution with
$\nu$ degrees of freedom, scaled to unit variance. Three models for the conditional variance
are compared. GARCH(1,1)~\cite{bollerslev1986generalized}:
\begin{equation}
\sigma_t^2 = \omega + \alpha\,\varepsilon_{t-1}^2 + \beta\,\sigma_{t-1}^2 .
\end{equation}
GJR-GARCH(1,1,1)~\cite{glosten1993relation}, where $I_{t-1}=1$ if $\varepsilon_{t-1}<0$ and
$0$ otherwise:
\begin{equation}
\sigma_t^2 = \omega + \left(\alpha + \gamma I_{t-1}\right)\varepsilon_{t-1}^2 + \beta\,\sigma_{t-1}^2 .
\end{equation}
EGARCH(1,1,1)~\cite{nelson1991conditional}:
\begin{equation}
\ln\sigma_t^2 = \omega + \alpha\left(|e_{t-1}| - \sqrt{2/\pi}\right) + \gamma\, e_{t-1} + \beta\ln\sigma_{t-1}^2 .
\end{equation}
In GJR a positive $\gamma$, and in EGARCH a negative $\gamma$, means that negative returns
raise volatility more than positive returns of the same size.

Parameters are estimated by maximum likelihood on the training block only. The fitted
parameters are then held fixed and the model is run forward over the whole series as a
one-step-ahead filter, so $\sigma_t$ on any validation or test day depends only on returns
up to day $t-1$ and on parameters estimated from training data. The standardized residual
is $z_t = (r_t - \mu)/\sigma_t$. All models are fitted with the \texttt{arch} Python package.

Model adequacy is checked on each block separately with Ljung-Box tests~\cite{ljung1978measure}
on $z_t$ and $z_t^2$ at 10 and 20 lags, an ARCH-LM test at 10 lags, and a Jarque-Bera test
of normality. The volatility model used for labels and inputs is chosen on the validation
diagnostics, not on training-sample BIC alone, because the preliminary results in
Chapter~4 show that the two can disagree. Before labels are fixed, the frozen-parameter
filter will be compared with rolling or expanding re-estimation.

\section{Anomaly Label Construction}

\subsection{Statistical Labels from Standardized Residuals}
A day is labelled anomalous if its standardized residual falls in the two-sided tail of
probability $p$ under the fitted unit-variance Student-$t$ distribution:
\begin{equation}
|z_t| > c, \qquad c = t_{\nu,\,1-p/2}\sqrt{\frac{\nu-2}{\nu}} ,
\end{equation}
where $t_{\nu,q}$ is the $q$ quantile of the standard Student-$t$ distribution. With
$p = 0.002$, a correctly specified model flags about one day in 500.

For comparison, a naive rule without a volatility model standardizes returns by the
training-block mean and standard deviation and flags days beyond the normal quantile for
the same tail probability.

\subsection{Other Label Sources (planned)}
\todo{Known market events: the event list and how its dates were sourced. Injected
anomalies: types, magnitudes and injection protocol. Freeze all three label rules in
docs/labeling\_rules.md before any detector is trained.}

\section{Detection Models}

\subsection{Baselines}
The first baselines are the two statistical rules above: a threshold on the raw-return
z-score and a threshold on the GARCH-standardized residual. \todo{Add classical detectors
if used.}

\subsection{Autoencoder on Raw versus GARCH-Standardized Returns (planned)}
\todo{Architecture, window length, training data, anomaly score $s_t$ and threshold
$\tau$. Both input types use identical settings, so the only difference is the input.}

\section{Transfer Learning Protocol (planned)}
\todo{Training on global indices, then applying to DSE data without retraining (zero-shot)
and after fine-tuning. State which parts are frozen and whether the volatility model is
re-estimated on DSE training data.}

\section{Evaluation Metrics}
Detectors are scored against labelled anomalous events. Precision, recall and $F_1$ are
\begin{equation}
Precision=\frac{TP}{TP+FP}, \qquad Recall=\frac{TP}{TP+FN},
\end{equation}
\begin{equation}
F_1=\frac{2\times Precision \times Recall}{Precision+Recall},
\end{equation}
where $TP$, $FP$ and $FN$ count true positives, false positives and false negatives.

Three versions of $F_1$ are reported together. Event-based $F_1$ counts each labelled event
once. Point-adjusted $F_1$ (F1-PA)~\cite{xu2018unsupervised} counts every point of an event as
detected if any one of its points is detected. PA\%K~\cite{kim2022towards} applies the
adjustment only when at least $K$\% of the event's points are detected. Because point
adjustment can make random scores look competitive~\cite{kim2022towards}, the gap between
F1-PA and the stricter scores is reported as a result in its own right.
\todo{Fix $K$ and decide whether to add VUS-PR~\cite{paparrizos2022volume}.}

\section{Reproducibility}
All settings, including split dates, model specifications and label thresholds, are kept in
one configuration file (\texttt{config.yaml}) in the project repository. Random seeds are
fixed at 42 for Python, NumPy and PyTorch. The code is public at
\url{https://github.com/smhasibul0/dse-anomaly-detection}. \todo{Software versions.}

\section{Ethical Considerations}
The study uses public market prices only and no personal data. The DSE website terms and
the licence of the Mendeley dataset have not been reviewed, so the code and download
scripts will be released but raw DSE prices will not be redistributed until the terms are
confirmed.

\section{Limitations of the Methodology}
The data are daily, so anomalies that begin and end within a trading day are invisible.
DSE prices are distorted by floor-price regimes and price limits, and the corporate-action
adjustment relies on an inferred rule. Statistical labels depend on the chosen volatility
model and threshold, and real manipulation cases on the DSE are not publicly labelled.
\todo{Add limitations found while running the experiments.}

\section{Summary of the Chapter}
\todo{Recap once the planned sections are complete.}
